\scanpage{60}

que \inserted{4 \deleted{\unsure{foncteurs}}\unsure{foncteurs} qui interviennent \unsure{au} \unsure{départ}}
les deux mbs de ces \deleted{appliqués} à $F$, sont
bien tous bornés \inserted{en \unsure{haut}}, ce qui ne semble
nullement clair. Mais par passage à la limite,
\inserted{par le th. de changement de base \unsure{connu}}
on peut supposer $X$ est de type fini sur $Z$.
\deleted{et \unsure{lisse} \unsure{sur} \unsure{affine} de la forme $\operatorname{Spec}Z[t_1,\unsure{\ldots|--},t_n]$.}

Alors le th de la dimension cohomologique
nous apprend que les
\deleted{$\unsure{\mathbb{R}}\,\unsure{cohomol}\,\unsure{\Gamma'}$}
\inserted{$(\mathbb{R}j_*)$}, $j$ ouvert
\unsure{Affinement} ! \deleted{\unsure{$\mathbb{R}$}}, \deleted{\unsure{ne} \unsure{dépend} \unsure{pas} \unsure{de} $\unsure{X^c}$ \unsure{peut-être}
\unsure{des} \unsure{raisons} pour \unsure{des} \unsure{coefficients} de $\unsure{2|l}$-torsion !!!}

! Ref exposé d'Artin.
\mspage{59}
